% ------------------------------------------------------------------
%  Capitolo 25 — Contro i miti
%  Parte VI
%  Ricerca di riferimento: ricerca 01 §14; 03 §6; 04 §8.4
%  Voci bibliografiche principali: pashler2008, dekker2012, bei2023, rayner2016, simons2016, letrud2018
%  Epigrafe: ✔ verificata sul testo (vedi ricerca 03 e registro, ricerca 00)
%  Scheda del capitolo: ricerca/06_indice-ragionato.md, capitolo 25
% ------------------------------------------------------------------
\chapter{Contro i miti}
\label{cap:miti}

\epigrafe[in ogni disciplina l'insegnamento della tecnica è debole senza la massima assiduità dell'esercizio]{in omni disciplina infirma est artis praeceptio sine summa adsiduitate exercitationis}{Rhetorica ad Herennium III, 24, 40}

\iniziale{Q}{uesto capitolo} \segnaposto{apertura: da scrivere — vedi la scheda nell'indice ragionato}.

\section{Gli stili di apprendimento}

\segnaposto{da scrivere}

\section{Il dieci per cento del cervello e altre leggende}

\segnaposto{da scrivere}

\section{La piramide dell'apprendimento}

\segnaposto{da scrivere}

\section{Lettura veloce e allenamento del cervello}

\segnaposto{da scrivere}

\section{Come riconoscere un metodo da quattro soldi}

\segnaposto{da scrivere}

\begin{daricordare}
\segnaposto{il principio del capitolo in due righe}
\end{daricordare}

\begin{domande}
  \item \segnaposto{domanda di recupero su questo capitolo}
  \item \segnaposto{domanda di applicazione a una materia che studi}
  \item \segnaposto{domanda di ripresa su un capitolo precedente}
\end{domande}

\begin{sintesi}
  \item \segnaposto{punto di sintesi}
\end{sintesi}

\finecapitolo
